% \begin{wrapfigure}{r}{0.5\textwidth}
% % \begin{figure}[t]
%     \centering
%     \includegraphics[width=\linewidth]{figures/discrimination_f.pdf}
%     \caption{\textbf{Qualitative visualization of learned attention and selected context/target regions using ViT-B/16 model.} 
%     % Rows for each sample denote the input image, attention map from the target encoder, context/target regions in I-JEPA and \ours. 
%     For \ours, the numbered regions (\Circled{1}–\Circled{5}) correspond to discriminative regions ordered by their estimated importance.
%     }
%     \label{fig:discrimination}
% % \end{figure}
% \end{wrapfigure}




\begin{figure}[t]
\centering
\begin{minipage}{0.57\linewidth}
  \centering
  \includegraphics[width=\linewidth]{figures/discrimination_f.pdf}
  \caption{\textbf{Qualitative visualization of learned attention and selected context/target regions using ViT-B/16 model.} 
    % Rows for each sample denote the input image, attention map from the target encoder, context/target regions in I-JEPA and \ours. 
    For \ours, the numbered regions (\Circled{1}–\Circled{5}) correspond to discriminative regions ordered by their estimated importance.}
  \label{fig:discrimination}
\end{minipage}\hfill
\begin{minipage}{0.4\linewidth}
  \centering
  \includegraphics[width=\linewidth]{figures/Patch_clustering.pdf}
  \caption{\textbf{Evolution of patch-level clustering during pre-training.} From left to right: input image and 4-cluster results from ViT-B/16 checkpoints at 150, 300, 450, and 600 epochs.}
  \label{fig:patch_cluster}
\end{minipage}
\end{figure}